\chapter{Complexity — Exp Poly}
%%% generated by the blueprint pipeline from TCSlib.Complexity.ClassNP.ExpPoly
%%%%%%%%%%%%%%%%%%%%%%%%%%%%%%%%%%%%%%%%%%%%%%%%%%%%%%%%%%%%%%%%%%%%%%%%%%%%%%%
\section{Overview}

This module introduces \emph{exponential-polynomial} bounds, i.e.\ functions
$T : \mathbb{N} \to \mathbb{N}$ with $T(n) \le 2^{K(n+1)^k}$ for some constants $K, k$.
It shows that such bounds contain the polynomials $a(n+1)^d$ and the exponentials
$c \cdot 2^{n^e}$, are closed under domination, sums, products and polynomial
reparametrization of the argument, and that every language decided within an
exponential-polynomial time bound lies in $\mathrm{EXP} = \bigcup_c \mathrm{DTIME}(2^{n^c})$
[AB09, Claim~2.4, \S2.6.2].

%%%%%%%%%%%%%%%%%%%%%%%%%%%%%%%%%%%%%%%%%%%%%%%%%%%%%%%%%%%%%%%%%%%%%%%%%%%%%%%
\section{Declarations}

\begin{definition}[Exponential-polynomial bound]
\lean{Complexity.ExpPoly}
\label{Complexity.ExpPoly}
\leanok
A function $T : \mathbb{N} \to \mathbb{N}$ is \emph{exponential-polynomial} if there exist
constants $K, k \in \mathbb{N}$ such that $T(n) \le 2^{K(n+1)^k}$ for every $n \in \mathbb{N}$.
\end{definition}

\begin{theorem}[Monotonicity of powers in the degree]
\lean{Complexity.pow_mono_deg}
\label{Complexity.pow_mono_deg}
\leanok
\difficulty{1}
For all natural numbers $n$, $k$ and $k'$ with $k \le k'$, we have $(n+1)^k \le (n+1)^{k'}$.
\end{theorem}

\begin{theorem}[Closure under domination]
\lean{Complexity.ExpPoly.of_le}
\label{Complexity.ExpPoly.of_le}
\leanok
\uses{Complexity.ExpPoly}
\difficulty{1}
Let $T, T' : \mathbb{N} \to \mathbb{N}$ satisfy $T(n) \le T'(n)$ for every $n$. If $T'$ is
exponential-polynomial, i.e.\ $T'(n) \le 2^{K(n+1)^k}$ for all $n$ and some constants
$K, k \in \mathbb{N}$, then $T$ is exponential-polynomial as well.
\end{theorem}

\begin{theorem}[Exponentials are exponential-polynomial]
\lean{Complexity.expPoly_exp}
\label{Complexity.expPoly_exp}
\leanok
\uses{Complexity.ExpPoly}
\difficulty{4}
For all constants $c, e \in \mathbb{N}$, the function $n \mapsto c \cdot 2^{n^e}$ is
exponential-polynomial: there exist constants $K, k \in \mathbb{N}$ such that
$c \cdot 2^{n^e} \le 2^{K(n+1)^k}$ for every $n \in \mathbb{N}$.
\end{theorem}

\begin{theorem}[Polynomials are exponential-polynomial]
\lean{Complexity.expPoly_poly}
\label{Complexity.expPoly_poly}
\leanok
\uses{Complexity.ExpPoly}
\difficulty{4}
For all constants $a, d \in \mathbb{N}$, the function $n \mapsto a(n+1)^d$ is
exponential-polynomial: there exist constants $K, k \in \mathbb{N}$ such that
$a(n+1)^d \le 2^{K(n+1)^k}$ for every $n \in \mathbb{N}$.
\end{theorem}

\begin{theorem}[Closure under sums]
\lean{Complexity.ExpPoly.add}
\label{Complexity.ExpPoly.add}
\leanok
\uses{Complexity.ExpPoly, Complexity.pow_mono_deg}
\difficulty{4}
Let $T_1, T_2 : \mathbb{N} \to \mathbb{N}$ both be exponential-polynomial, i.e.\ each is
bounded by $2^{K(n+1)^k}$ for all $n$ and some constants $K, k \in \mathbb{N}$. Then the sum
$n \mapsto T_1(n) + T_2(n)$ is exponential-polynomial.
\end{theorem}

\begin{theorem}[Closure under products]
\lean{Complexity.ExpPoly.mul}
\label{Complexity.ExpPoly.mul}
\leanok
\uses{Complexity.ExpPoly, Complexity.pow_mono_deg}
\difficulty{4}
Let $T_1, T_2 : \mathbb{N} \to \mathbb{N}$ both be exponential-polynomial, i.e.\ each is
bounded by $2^{K(n+1)^k}$ for all $n$ and some constants $K, k \in \mathbb{N}$. Then the
product $n \mapsto T_1(n) \cdot T_2(n)$ is exponential-polynomial.
\end{theorem}

\begin{theorem}[Composition with a polynomial argument]
\lean{Complexity.ExpPoly.comp_poly}
\label{Complexity.ExpPoly.comp_poly}
\leanok
\uses{Complexity.ExpPoly}
\difficulty{4}
Let $T : \mathbb{N} \to \mathbb{N}$ be exponential-polynomial, i.e.\ $T(n) \le 2^{K(n+1)^k}$
for all $n$ and some constants $K, k \in \mathbb{N}$. Then for all constants
$a, d \in \mathbb{N}$, the function $n \mapsto T\bigl(a(n+1)^d\bigr)$ is also
exponential-polynomial.
\end{theorem}

\begin{theorem}[Exponential-polynomial time lies in $\mathrm{EXP}$]
\lean{Complexity.ExpPoly.mem_EXP}
\label{Complexity.ExpPoly.mem_EXP}
\leanok
\uses{Complexity.DTIME, Complexity.EXP, Complexity.ExpPoly, Turing.FinTM.ComputesInTime.mono}
\difficulty{5}
Let $T : \mathbb{N} \to \mathbb{N}$ be exponential-polynomial, i.e.\ $T(n) \le 2^{K(n+1)^k}$
for all $n$ and some constants $K, k \in \mathbb{N}$, and let $L$ be a language over the
binary alphabet in $\mathrm{DTIME}(T)$, i.e.\ decided by some finite multi-tape Turing machine
within $c \cdot T(n)$ steps on inputs of length $n$, for some constant $c$. Then
$L \in \mathrm{EXP}$, that is, $L \in \mathrm{DTIME}(2^{n^{c'}})$ for some constant
$c' \in \mathbb{N}$.
\end{theorem}
